\section{Introduction}
\label{sec:1}

A strongly coupled plasma with a non-abelian global symmetry can be
magnetically unstable. In the simplest holographic model, Einstein--$SU(2)$
Yang--Mills theory in $\mathrm{AdS}_5$, a magnetic field in the abelian
direction of the gauge group makes the charged vector bosons tachyonic in
their lowest Landau level. This is the holographic form of the
Nielsen--Olesen instability~\cite{Nielsen:1978rm}. In the probe limit, where
the gauge field does not deform the geometry, the plasma condenses once the
field exceeds $B_c \simeq 50\,T^2$~\cite{Ammon:2011je}. The condensate
carries the flux in a lattice of vortices, like the Abrikosov lattice of a
type II superconductor~\cite{Abrikosov:1956sx,Kleiner:1964zz}, and numerical
comparisons of candidate lattices found the triangular one
lowest~\cite{Bu:2012mq}. We show in appendix~\ref{app:probe} that this is a
theorem: the quartic vortex interaction of the probe theory is completely
monotone in the momentum transfer, and Montgomery's
theorem~\cite{Montgomery:1988mt,Cohn:2007universally,Betermin:2016theta}
then makes the triangular lattice the minimum among all one-flux-quantum
Bravais lattices.

All of this ignores the gravity of the field itself. The field sources the
metric with a strength set by the single coupling of the theory,
$\alpha = \kappa^2/\hat g^2$, and enters the geometry through
$\beta = \alpha B^2$. Since $B_c \simeq 5$ in horizon units, $\beta$ is of
order one already at couplings of a few percent: $\beta = 1$ at the onset
when $\alpha = 0.028$. This paper asks what happens once the field bends
the geometry. The answer is that gravity can remove the linear instability
altogether.

The calculation is exact in the coupling and expands only in the
condensate, which is small near the onset. This is possible because the
condensate-free background at any $\beta$ is known: it is the magnetic
brane of D'Hoker and Kraus~\cite{DHoker:2009mmn,DHoker:2009ixq}. The onset
becomes a Sturm--Liouville problem on the brane, and the quartic term a
boundary value problem for the brane's linear response.

The central result is that the critical field rises with the coupling and
diverges at
\begin{equation}
\alpha_\star = \tfrac23~,
\end{equation}
above which the normal phase of the magnetised plasma has no linear
instability in the charged sector, at any field or temperature. The reason
is simple. At zero temperature the interior of the brane is an
$\mathrm{AdS}_3 \times \mathbb{R}^2$ throat in which the flux plane
stretches in proportion to $B$~\cite{DHoker:2009mmn}. The field measured
in throat units, and with it the tachyonic mass of the lowest Landau level,
is therefore a number fixed by $\alpha$, and the onset condition becomes
the Breitenlohner--Freedman
bound~\cite{Breitenlohner:1982bm,Breitenlohner:1982jf} of the
$\mathrm{AdS}_3$ factor. The
tachyon violates it exactly when $\alpha < 2/3$.

Violating the bound is sufficient for the instability; a node count
shows that here it is also necessary, and a positivity argument extends
the result to the whole charged sector. At finite temperature we establish the result
numerically. In $d + 1$ bulk dimensions the bound is
$\alpha_\star(d) = 16/(d(d-1)(d-2))$, and it is the critical coupling for
$d = 4, 5, 6$. Since $\alpha$ is proportional to the ratio of the current
and stress-tensor central charges of the boundary theory, the criterion can
be stated in boundary terms: for $d = 4, 5, 6$ a magnetised plasma with an
Einstein--Yang--Mills dual has a charged instability at zero temperature
if and only if
\begin{equation}
\frac{C_J}{C_T} < \frac{8}{d^2(d+1)}~,
\end{equation}
which in four dimensions says that the current carries less than one tenth
of the stress-tensor central charge. The onset temperature vanishes at
$\alpha_\star$ with an essential singularity of
Berezinskii--Kosterlitz--Thouless type, whose coefficient we derive, though it sets in only at extremely low
temperatures.

Below $\alpha_\star$ gravity adds an attractive channel to the interaction
between vortices: the exchange of the plasma's stress tensor between
density modulations of the condensate. The attraction barely changes the
shape of the lattice, but it decides its stability. At
\begin{equation}
\alpha_L = 0.372
\end{equation}
the quartic coefficient of the uniform lattice changes sign, the triangular
lattice becomes unstable to long-wavelength density modulation, and the
transition at $B_c$ turns first order, so $\alpha_L$ is a tricritical
point. The plasma thus passes through three
regimes (figure~\ref{fig:1}): a vortex crystal below $\alpha_L$, a
first-order regime up to $\alpha_\star$, in which $B_c$ is only a spinodal,
and no charged instability beyond. Whether a condensed phase survives above
$\alpha_\star$, reached by a first-order jump, is the main question we
leave open.

\begin{figure}[t]
\centering
\includegraphics{poster_bc_alpha}
\caption{The critical field $B_c(\alpha)$ at fixed temperature, on a
logarithmic axis: the collocation ladder (circles) and, beyond its top rung
at $\alpha = 0.591$, direct solves by shooting from the horizon (grey). The
curve rises convexly and has a vertical asymptote at the critical coupling
$\astar = 2/3$. The uniform vortex crystal forms in a second-order
transition for $\alpha < \alpha_L$; for $\alpha_L < \alpha < \astar$ the
transition is first order and $B_c$ is a spinodal; for $\alpha > \astar$
the normal phase has no linear charged instability. The right axis reads
the same curve as the onset temperature at fixed field,
$T_c/\sqrt B = 1/(\pi\sqrt{B_c u_H^2})$, which falls from $0.141$ in the
probe limit to zero at $\astar$.}
\label{fig:1}
\end{figure}

Einstein--$SU(2)$ Yang--Mills holography goes back
to~\cite{Gubser:2008zu,Gubser:2008wv}, where a chemical potential drives a
$p$-wave condensate; in the backreacted model of~\cite{Gubser:2008zu} the
critical temperature was found numerically to vanish at a critical
coupling. The magnetic brane and its throat are due to D'Hoker and
Kraus~\cite{DHoker:2009mmn,DHoker:2012rlj}, who computed its correlators
with Shah~\cite{DHoker:2010xwl} and also found a
quantum critical point on the charged brane, tuned by field over charge
density at fixed Chern--Simons coupling~\cite{DHoker:2010zpp,DHoker:2010onp}.
The stability of the throat within $D = 5$ $SO(6)$ gauged supergravity,
which uplifts to type IIB on $S^5$, was examined
in~\cite{Almuhairi:2010rb,Almuhairi:2011ws,Donos:2011pn}. For the
equal-charge embedding of D'Hoker and Kraus the scalars of the truncation
are stable~\cite{Almuhairi:2010rb}, but the $SO(6)$ $W$ bosons have a
Nielsen--Olesen mode that violates the $\mathrm{AdS}_3$
bound~\cite{Donos:2011pn},\footnote{The mode was first found in v1
of~\cite{Almuhairi:2011ws}.} and, including their mixing with
the charged scalars, \cite{Donos:2011pn} mapped the unstable embeddings and
found only a narrow window of non-supersymmetric throats with no known
instability next to the supersymmetric ones. In all of these
the coupling is fixed by supergravity and the $U(1)$ embedding is the
parameter. Here the embedding is fixed and the coupling is the parameter,
which is what produces a critical value. In $d = 3$,
\cite{Wong:2013rda} derived the corresponding bound on the magnetic
$\mathrm{AdS}_4$ brane, which our general-$d$ formula reproduces, computed
the onset at finite coupling and temperature, and studied the lattice in
the probe limit. In $d = 3$ we find that the bound is not the critical
coupling (section~\ref{sec:4.3}): a bound state outside the $\mathrm{AdS}_2$
throat keeps the plasma unstable up to $\alpha = 3.037$ rather than $8/3$.

Backreacted vortex lattices of abelian condensates have been constructed
near onset~\cite{Bao:2013fda}, in a circular-cell
approximation~\cite{Tallarita:2019amp} and in full~\cite{Donos:2020viz},
which finds the triangular lattice preferred; the onset on a backreacted magnetic brane
has been studied for an abelian $s$-wave condensate~\cite{Chernicoff:2024ymn},
and a single backreacted vortex with non-abelian orientational moduli was
constructed in~\cite{Tallarita:2019czh}. With an adjoint Higgs field the
$SU(2)$ theory on the magnetic $\mathrm{AdS}_4$ brane instead forms a
monopole wall through a first-order transition~\cite{Bolognesi:2010nb}. The
flat-space vortex lattices of the gluon~\cite{Ambjorn:1980ms} and
electroweak $W$~\cite{Ambjorn:1988fx,Ambjorn:1988tm}, and the proposed
$\rho$-meson condensate of QCD~\cite{Chernodub:2010qx,Adhikari:2024bfa},
are the same mechanism without gravity; with gravity, the same instability
appears near the horizon of a magnetically charged black
hole~\cite{Lee:1991qs,Maldacena:2020skw}. The $\rho$-meson onset has also
been found holographically in the Sakai--Sugimoto
model~\cite{Callebaut:2011ab,Callebaut:2013wba}. Earlier holographic vortex lattices in the probe
limit are~\cite{Maeda:2009vf,Murray:2011gr}. The holographic
Berezinskii--Kosterlitz--Thouless scaling we find is that
of~\cite{Kaplan:2009kr,Jensen:2010ga,Iqbal:2010eh}, with a bulk coupling
tuned through the bound rather than a filling fraction, an
operator dimension, a field~\cite{Faulkner:2010gj} or a flavour
number~\cite{Alho:2013dka}.

Section~\ref{sec:2} sets up the theory and the expansion,
section~\ref{sec:3} gives the onset curve $B_c(\alpha)$,
section~\ref{sec:4} derives the critical coupling, and section~\ref{sec:5}
treats the lattice. Section~\ref{sec:6} lists open questions.
